\subsection{Семинары «Порождающие паттерны» (4 очных часа)}

\begin{itemize}
   \item Используя паттерны <<Фабричный метод>> и <<Абстрактная фабрика>> 
реализуйте абстрактное создание и работу с различными сущностями в вашей программе 
(в каждом варианте по две опции сущностей). 
По результатам использования паттернов у вас должно упроститься добавление новых видов
сущностей
   \item Используя паттерны <<строитель>> и <<команда>> (этот паттерн мы используем
   заранее) реализуйте создание меню из любого количества пунктов с любыми выполняемыми
   действиями
   \item Реализуйте запись и чтение из файла вашей базы данных (формат JSON)
   в следующей логике: 
   после открытия программы пользователь выбирает: создать новую БД (и имя файла), 
   открыть старую БД (и имя файла), работать в старом режиме (без файла). 
   Файл открывается только тогда, когда он первый раз нужен и дальше не закрывается, 
   пока пользователь явно не выберет операцию закрытия файла. Здесь используется паттерн
   <<ленивая инициализация>>. 
   \item Добавить возможность клонировать сущность (создается копия и можно отредактировать
   поля по выбору: проще дать возможность редактировать все поля, и если пользователь 
   просто нажмет Enter, то использовать старое значение). Здесь мы используем паттерн
   <<Прототип>>.
\end{itemize}